\documentclass[tikz]{standalone}
\usepackage{amsmath}
\definecolor{bgcol}{HTML}{171D2B}
\definecolor{accent}{HTML}{F7B538}
\definecolor{setcol}{HTML}{7DA2C9}
\pagecolor{bgcol}
\begin{document}
\begin{tikzpicture}[scale=1.0, color=white, font=\large]
  % Dua garis berpotongan: sudut bertolak belakang 50 dan 130 derajat
  % Garis AB mendatar, garis CD miring 50 derajat (koordinat polar)
  % Sudut kanan atas 50 derajat, kiri bawah 50 derajat (busur tunggal)
  % Sudut kiri atas dan kanan bawah 130 derajat (busur ganda)
  \coordinate (O) at (0, 0);

  % Kedua garis
  \draw[thick] (-3.5, 0) -- (3.5, 0);
  \draw[thick] (230:3.5) -- (50:3.5);

  % Busur TUNGGAL untuk pasangan sudut 50 derajat
  \draw (O) ++(0:0.8)   arc (0:50:0.8);
  \draw (O) ++(180:0.8) arc (180:230:0.8);
  % Busur GANDA untuk pasangan sudut 130 derajat
  \draw (O) ++(50:1.05)  arc (50:180:1.05);
  \draw (O) ++(50:1.25)  arc (50:180:1.25);
  \draw (O) ++(230:1.05) arc (230:360:1.05);
  \draw (O) ++(230:1.25) arc (230:360:1.25);

  % Titik potong
  \fill (O) circle (2pt);
  \node[below right] at (O) {$O$};

  % Label ujung-ujung garis
  \node[left]  at (-3.5, 0) {$A$};
  \node[right] at (3.5, 0)  {$B$};
  \node[above] at (50:3.8)  {$C$};
  \node[below] at (230:3.8) {$D$};

  % Label sudut: 50 derajat (busur tunggal), 130 derajat (busur ganda)
  \node at (25:1.2)  {$50^{\circ}$};
  \node at (205:1.2) {$50^{\circ}$};
  \node at (115:1.6) {$130^{\circ}$};
  \node at (295:1.6) {$130^{\circ}$};
\end{tikzpicture}
\end{document}
